\section{Homework Midterm}

\begin{exercise}
    Let \(R\) be a ring with \(1\).
    
    \begin{enumerate}
        \item
        \begin{enumerate}
            \item[(i)]
            Let
            \[
                h:U\longrightarrow V
            \]
            be a monomorphism of \(R\)-modules. Then the following two conditions are equivalent.
            \begin{enumerate}
                \item[(a)]
                A homomorphism
                \[
                    k:V\longrightarrow Y
                \]
                must be a monomorphism whenever the composite map
                \[
                    U \xrightarrow{\,h\,} V \xrightarrow{\,k\,} Y
                \]
                is a monomorphism.
    
                \item[(b)]
                If \(0\neq W\subset V\), then
                \[
                    h(U)\cap W\neq 0.
                \]
            \end{enumerate}
    
            The monomorphism \(h\) satisfying the above conditions is said to be
            \emph{essential}. If there is an injective \(R\)-module \(Q\) with an essential
            monomorphism
            \[
                h:V\longrightarrow Q,
            \]
            then the diagram
            \[
                V \xrightarrow{\,h\,} Q
            \]
            is called an \emph{injective hull} of \(V\).
    
            \item[(ii)]
            Prove that any \(R\)-module \(V\) has an injective hull.
        \end{enumerate}
    
        \item
        \begin{enumerate}
            \item[(i)]
            Prove that the following holds for an idempotent \(e\in R\):
            \[
                J(eRe)=eJ(R)e=eRe\cap J(R).
            \]
    
            \item[(ii)]
            Prove that
            \[
                J(M_n(R))=M_n(J(R)).
            \]
            Here \(M_n(R)\) means the matrix ring.
        \end{enumerate}
    
        \item
        Prove that a submodule of a finitely generated free module over a PID is free.
    \end{enumerate}
    \end{exercise}

    \begin{proof}[Proof of 1(i)]
        We identify \(U\) with its image \(h(U)\subseteq V\).
        
        First assume (a). Suppose that there exists a nonzero submodule
        \(W\subseteq V\) such that
        \[
            h(U)\cap W=0.
        \]
        Let
        \[
            q:V\longrightarrow V/W
        \]
        be the canonical quotient map. Since \(h(U)\cap W=0\), the composite
        \[
            U\xrightarrow{\,h\,}V\xrightarrow{\,q\,}V/W
        \]
        is a monomorphism. Hence by (a), \(q\) must be a monomorphism. But
        \[
            \ker q=W\neq 0,
        \]
        a contradiction. Therefore for every nonzero submodule \(W\subseteq V\), one has
        \[
            h(U)\cap W\neq 0.
        \]
        Thus (b) holds.
        
        Conversely, assume (b). Let
        \[
            k:V\longrightarrow Y
        \]
        be an \(R\)-homomorphism such that \(kh\) is a monomorphism. We prove that \(k\)
        is a monomorphism.
        
        Let
        \[
            W=\ker k.
        \]
        If \(W\neq 0\), then by (b),
        \[
            h(U)\cap W\neq 0.
        \]
        Choose \(0\neq h(u)\in h(U)\cap W\). Then
        \[
            k(h(u))=0.
        \]
        Since \(kh\) is a monomorphism, this implies \(u=0\), hence \(h(u)=0\), a contradiction.
        Therefore \(W=0\), so \(k\) is a monomorphism.
        
        Hence (a) and (b) are equivalent.
        \end{proof}
        
        \begin{definition}[Essential extension]
            Let \(V\subseteq E\) be an inclusion of \(R\)-modules. We say that \(V\) is
            \emph{essential} in \(E\), or that \(E\) is an \emph{essential extension} of \(V\),
            if for every nonzero submodule \(W\subseteq E\), one has
            \[
                W\cap V\neq 0.
            \]
            Equivalently, every nonzero submodule of \(E\) must meet \(V\) nontrivially.
            \end{definition}
            
            \begin{remark}
            Intuitively, the condition \(V\subseteq E\) being essential means that \(E\) contains
            no nonzero submodule which is completely disjoint from \(V\). Thus \(E\) is obtained
            from \(V\) in a way that does not create an independent ``new part'' away from \(V\).
            \end{remark}
            
            \begin{lemma}[Existence of a maximal essential extension inside a fixed ambient module]
            Let \(V\) be an \(R\)-module, and suppose that \(V\) is embedded into an \(R\)-module
            \(I\):
            \[
                V\subseteq I.
            \]
            Consider the set
            \[
                \mathcal S
                =
                \{\,E\mid V\subseteq E\subseteq I,\ \text{and }V\subseteq E
                \text{ is an essential extension}\,\}.
            \]
            Then \(\mathcal S\) has a maximal element.
            \end{lemma}
            
            \begin{proof}
            First, \(\mathcal S\neq \varnothing\). Indeed, \(V\in \mathcal S\), because
            \(V\subseteq V\) is clearly essential: if \(0\neq W\subseteq V\), then
            \[
                W\cap V=W\neq 0.
            \]
            
            We partially order \(\mathcal S\) by inclusion. To apply Zorn's lemma, we show that
            every chain in \(\mathcal S\) has an upper bound in \(\mathcal S\).
            
            Let
            \[
                \{E_\alpha\}_{\alpha\in A}
            \]
            be a chain in \(\mathcal S\). This means that for any \(\alpha,\beta\in A\), either
            \[
                E_\alpha\subseteq E_\beta
                \quad\text{or}\quad
                E_\beta\subseteq E_\alpha.
            \]
            Define
            \[
                E=\bigcup_{\alpha\in A}E_\alpha.
            \]
            Because the family \(\{E_\alpha\}_{\alpha\in A}\) is a chain, the union \(E\) is again
            an \(R\)-submodule of \(I\). Clearly,
            \[
                V\subseteq E\subseteq I.
            \]
            
            It remains to prove that \(V\subseteq E\) is essential. Let
            \[
                0\neq W\subseteq E
            \]
            be a nonzero submodule. Choose a nonzero element
            \[
                0\neq w\in W.
            \]
            Since \(E=\bigcup_{\alpha\in A}E_\alpha\), there exists some \(\alpha\in A\) such that
            \[
                w\in E_\alpha.
            \]
            Hence
            \[
                Rw\subseteq E_\alpha.
            \]
            Because \(w\neq 0\) and \(R\) has identity, we have
            \[
                Rw\neq 0.
            \]
            Since \(E_\alpha\in \mathcal S\), the inclusion \(V\subseteq E_\alpha\) is essential.
            Therefore
            \[
                Rw\cap V\neq 0.
            \]
            But \(Rw\subseteq W\), so
            \[
                W\cap V\neq 0.
            \]
            Thus \(V\subseteq E\) is essential, and hence
            \[
                E\in \mathcal S.
            \]
            
            Therefore every chain in \(\mathcal S\) has an upper bound in \(\mathcal S\). By Zorn's
            lemma, \(\mathcal S\) has a maximal element. Denote it by \(E\).
            \end{proof}
            
            \begin{remark}
            The maximality of \(E\) means the following: if \(E'\) is another submodule of \(I\)
            such that
            \[
                E\subseteq E'\subseteq I
            \]
            and \(V\subseteq E'\) is still essential, then necessarily
            \[
                E'=E.
            \]
            Equivalently, there is no strictly larger submodule \(E'\) with
            \[
                E\subsetneq E'\subseteq I
            \]
            such that \(V\subseteq E'\) remains an essential extension.
            \end{remark}
            
            \begin{remark}
            In the proof of the existence of injective hulls, one usually first uses the standard
            fact that every \(R\)-module \(V\) can be embedded into some injective module \(I\).
            Then the above lemma gives a maximal essential extension
            \[
                V\subseteq E\subseteq I.
            \]
            The remaining step is to prove that this maximal essential extension \(E\) is itself
            injective. Once this is done, the inclusion
            \[
                V\hookrightarrow E
            \]
            is an injective hull of \(V\).
            \end{remark}

            \begin{lemma}[Baer's criterion]
                Let \(R\) be a ring with identity, and let \(E\) be a left \(R\)-module. Then \(E\)
                is injective if and only if for every left ideal \(\mathfrak a\subseteq R\), every
                \(R\)-homomorphism
                \[
                    f:\mathfrak a\longrightarrow E
                \]
                can be extended to an \(R\)-homomorphism
                \[
                    \widetilde f:R\longrightarrow E.
                \]
                Equivalently, for every left ideal \(\mathfrak a\subseteq R\), every diagram
                \[
                \begin{tikzcd}
                \mathfrak a \arrow[r,"f"] \arrow[d,hook] & E \\
                R \arrow[ru,dashed,"\widetilde f"'] &
                \end{tikzcd}
                \]
                can be completed by some \(\widetilde f:R\to E\).
                \end{lemma}
                
                \begin{proof}
                First suppose that \(E\) is injective. Let \(\mathfrak a\subseteq R\) be a left ideal,
                and let
                \[
                    f:\mathfrak a\longrightarrow E
                \]
                be an \(R\)-homomorphism. Since \(\mathfrak a\hookrightarrow R\) is a monomorphism
                of left \(R\)-modules, the injectivity of \(E\) implies that \(f\) extends to an
                \(R\)-homomorphism
                \[
                    \widetilde f:R\longrightarrow E.
                \]
                Thus the stated condition holds.
                
                Conversely, assume that every homomorphism from a left ideal of \(R\) to \(E\) can be
                extended to \(R\). We prove that \(E\) is injective.
                
                Let \(A\subseteq B\) be an inclusion of left \(R\)-modules, and let
                \[
                    f:A\longrightarrow E
                \]
                be an \(R\)-homomorphism. We will prove that \(f\) extends to \(B\).
                
                Consider the set \(\mathcal S\) of pairs \((C,g)\), where
                \[
                    A\subseteq C\subseteq B
                \]
                is a submodule and
                \[
                    g:C\longrightarrow E
                \]
                is an \(R\)-homomorphism satisfying
                \[
                    g|_A=f.
                \]
                We partially order \(\mathcal S\) by declaring
                \[
                    (C,g)\leq (C',g')
                \]
                if
                \[
                    C\subseteq C'
                    \quad\text{and}\quad
                    g'|_C=g.
                \]
                
                The set \(\mathcal S\) is nonempty, since \((A,f)\in \mathcal S\). If
                \[
                    \{(C_i,g_i)\}_{i\in I}
                \]
                is a chain in \(\mathcal S\), then
                \[
                    C=\bigcup_{i\in I}C_i
                \]
                is a submodule of \(B\), and we may define
                \[
                    g:C\longrightarrow E
                \]
                by
                \[
                    g(c)=g_i(c)
                \]
                whenever \(c\in C_i\). This is well-defined because the pairs form a chain. Hence
                \((C,g)\) is an upper bound of the chain. By Zorn's lemma, \(\mathcal S\) has a
                maximal element, say
                \[
                    (C,g).
                \]
                
                We claim that \(C=B\). Suppose not. Choose
                \[
                    b\in B\setminus C.
                \]
                Define
                \[
                    \mathfrak a=\{\,r\in R\mid rb\in C\,\}.
                \]
                Then \(\mathfrak a\) is a left ideal of \(R\). Indeed, if \(r,s\in\mathfrak a\), then
                \[
                    (r+s)b=rb+sb\in C,
                \]
                and if \(t\in R\), then
                \[
                    (tr)b=t(rb)\in C.
                \]
                
                Define an \(R\)-homomorphism
                \[
                    h:\mathfrak a\longrightarrow E
                \]
                by
                \[
                    h(r)=g(rb).
                \]
                This is well-defined and \(R\)-linear because \(g\) is \(R\)-linear.
                
                By the assumed condition, \(h\) extends to an \(R\)-homomorphism
                \[
                    \widetilde h:R\longrightarrow E.
                \]
                Let
                \[
                    e=\widetilde h(1).
                \]
                Then for every \(r\in R\),
                \[
                    \widetilde h(r)=re,
                \]
                and in particular, for every \(r\in\mathfrak a\),
                \[
                    g(rb)=h(r)=\widetilde h(r)=re.
                \]
                
                Now define
                \[
                    g':C+Rb\longrightarrow E
                \]
                by
                \[
                    g'(c+rb)=g(c)+re,
                    \qquad c\in C,\ r\in R.
                \]
                We check that \(g'\) is well-defined. Suppose
                \[
                    c+rb=c'+r'b.
                \]
                Then
                \[
                    (r-r')b=c'-c\in C,
                \]
                so \(r-r'\in\mathfrak a\). Hence
                \[
                    g((r-r')b)=(r-r')e.
                \]
                But
                \[
                    g((r-r')b)=g(c'-c)=g(c')-g(c).
                \]
                Therefore
                \[
                    g(c')-g(c)=(r-r')e,
                \]
                which implies
                \[
                    g(c)+re=g(c')+r'e.
                \]
                Thus \(g'\) is well-defined. It is straightforward to check that \(g'\) is an
                \(R\)-homomorphism.
                
                Moreover, \(g'\) extends \(g\), since
                \[
                    g'(c)=g'(c+0b)=g(c).
                \]
                Thus
                \[
                    (C,g)<(C+Rb,g')
                \]
                unless \(C+Rb=C\). But \(b\notin C\), so \(C+Rb\neq C\). This contradicts the
                maximality of \((C,g)\).
                
                Therefore \(C=B\). Hence \(f:A\to E\) extends to a homomorphism \(B\to E\).
                This proves that \(E\) is injective.
                \end{proof}

        \begin{proof}[Proof of 1(ii)]
        
        It remains to prove that \(E\) is injective. We use Baer's criterion. Suppose, for
        contradiction, that \(E\) is not injective. Then there exists a left ideal
        \[
            I\subseteq R
        \]
        and an \(R\)-homomorphism
        \[
            f:I\longrightarrow E
        \]
        which cannot be extended to a homomorphism \(R\to E\).
        
        Construct an \(R\)-module
        \[
            P=(E\oplus R)/N,
        \]
        where
        \[
            N=\{(f(a),-a)\mid a\in I\}.
        \]
        The natural map
        \[
            E\longrightarrow P,\qquad e\longmapsto (e,0)+N
        \]
        is injective. Indeed, if \((e,0)\in N\), then
        \[
            (e,0)=(f(a),-a)
        \]
        for some \(a\in I\), hence \(a=0\) and \(e=f(0)=0\).
        
        Choose, by Zorn's lemma, a submodule \(C\subseteq P\) maximal subject to
        \[
            C\cap E=0.
        \]
        Then the image of \(E\) in \(P/C\) is essential. Indeed, if \(D/C\) is a nonzero
        submodule of \(P/C\), then \(D\supsetneq C\). By maximality of \(C\), we must have
        \[
            D\cap E\neq 0.
        \]
        Hence
        \[
            (D/C)\cap E\neq 0.
        \]
        Thus \(E\subseteq P/C\) is an essential extension.
        
        Since \(V\subseteq E\) is essential and \(E\subseteq P/C\) is essential, by transitivity
        \(V\subseteq P/C\) is essential. But \(E\) was chosen to be a maximal essential
        extension of \(V\). Therefore
        \[
            P/C=E.
        \]
        
        Let
        \[
            x=(0,1)+N+C\in P/C.
        \]
        Since \(P/C=E\), there exists \(e\in E\) such that
        \[
            x=e.
        \]
        For every \(a\in I\), we have in \(P/C\)
        \[
            ax=(0,a)+N+C=(f(a),0)+N+C=f(a).
        \]
        On the other hand, since \(x=e\), we also have
        \[
            ax=ae.
        \]
        Therefore
        \[
            f(a)=ae
            \qquad
            \text{for all }a\in I.
        \]
        Hence \(f\) extends to an \(R\)-homomorphism
        \[
            R\longrightarrow E,\qquad r\longmapsto re,
        \]
        contradicting the choice of \(f\).
        
        Thus \(E\) is injective. Since \(V\subseteq E\) is an essential monomorphism and
        \(E\) is injective, the inclusion
        \[
            V\hookrightarrow E
        \]
        is an injective hull of \(V\).
        \end{proof}
        
        \begin{proof}[Proof of 2(i)]
        We prove
        \[
            J(eRe)=eJ(R)e=eRe\cap J(R).
        \]
        
        First observe that
        \[
            eJ(R)e\subseteq eRe\cap J(R),
        \]
        because \(J(R)\) is a two-sided ideal of \(R\). Conversely, if
        \[
            x\in eRe\cap J(R),
        \]
        then \(x=exe\), so
        \[
            x\in eJ(R)e.
        \]
        Hence
        \[
            eJ(R)e=eRe\cap J(R).
        \]
        
        It remains to prove
        \[
            J(eRe)=eJ(R)e.
        \]
        
        First let \(a\in eJ(R)e\). Then \(a\in eRe\). For every \(y\in eRe\), we have
        \[
            ya\in J(R).
        \]
        By the usual characterization of the Jacobson radical, \(1-ya\) has a left inverse in
        \(R\). Thus there exists \(u\in R\) such that
        \[
            u(1-ya)=1.
        \]
        Multiplying by \(e\) on both sides gives
        \[
            eue(e-ya)=e.
        \]
        Hence \(e-ya\) has a left inverse in the ring \(eRe\), whose identity is \(e\). Therefore
        \[
            a\in J(eRe).
        \]
        Thus
        \[
            eJ(R)e\subseteq J(eRe).
        \]
        
        Conversely, let
        \[
            a\in J(eRe).
        \]
        We show that \(a\in J(R)\). It is enough to show that \(a\) annihilates every simple
        left \(R\)-module.
        
        Let \(S\) be a simple left \(R\)-module. Then \(eS\) is either \(0\) or a simple
        left \(eRe\)-module. Indeed, if \(0\neq T\subseteq eS\) is an \(eRe\)-submodule and
        \(0\neq t\in T\), then \(Rt=S\), and hence
        \[
            eS=eRt=eRe\,t\subseteq T.
        \]
        Therefore \(T=eS\).
        
        Since \(a\in J(eRe)\), it annihilates every simple \(eRe\)-module. Hence
        \[
            a(eS)=0.
        \]
        But \(a\in eRe\), so \(a=ae\). Therefore, for every \(s\in S\),
        \[
            as=ae s=a(es)=0.
        \]
        Thus \(a\) annihilates every simple left \(R\)-module, and hence
        \[
            a\in J(R).
        \]
        Since also \(a\in eRe\), we get
        \[
            a\in eRe\cap J(R)=eJ(R)e.
        \]
        Therefore
        \[
            J(eRe)\subseteq eJ(R)e.
        \]
        
        Combining the two inclusions gives
        \[
            J(eRe)=eJ(R)e=eRe\cap J(R).
        \]
        \end{proof}
        
        \begin{proof}[Proof of 2(ii)]
        Let
        \[
            S=M_n(R).
        \]
        We prove
        \[
            J(S)=M_n(J(R)).
        \]
        
        Let \(E_{ij}\) denote the standard matrix units. Applying part \(2(i)\) to the
        idempotent \(E_{11}\in S\), we obtain
        \[
            J(E_{11}SE_{11})=E_{11}J(S)E_{11}.
        \]
        But
        \[
            E_{11}SE_{11}\cong R,
        \]
        so
        \[
            E_{11}J(S)E_{11}\cong J(R).
        \]
        
        Now take \(A=(a_{ij})\in J(S)\). Then
        \[
            E_{1i}AE_{j1}\in E_{11}J(S)E_{11}.
        \]
        But
        \[
            E_{1i}AE_{j1}=a_{ij}E_{11}.
        \]
        Therefore \(a_{ij}\in J(R)\) for all \(i,j\). Hence
        \[
            J(S)\subseteq M_n(J(R)).
        \]
        
        Conversely, we show that
        \[
            M_n(J(R))\subseteq J(S).
        \]
        It is enough to prove that for every \(a\in J(R)\) and every \(i,j\),
        \[
            aE_{ij}\in J(S).
        \]
        
        Let \(X=(x_{pq})\in S\). We must show that
        \[
            I_n-X(aE_{ij})
        \]
        has a left inverse in \(S\).
        
        Set
        \[
            u=
            \begin{pmatrix}
                x_{1i}a\\
                x_{2i}a\\
                \vdots\\
                x_{ni}a
            \end{pmatrix},
            \qquad
            v^{T}=
            \begin{pmatrix}
                0&\cdots&0&1&0&\cdots&0
            \end{pmatrix},
        \]
        where the \(1\) occurs in the \(j\)-th position. Then
        \[
            X(aE_{ij})=uv^{T},
        \]
        and
        \[
            v^{T}u=x_{ji}a.
        \]
        Since \(a\in J(R)\), the element
        \[
            1-x_{ji}a
        \]
        has a left inverse in \(R\). Hence there exists \(c\in R\) such that
        \[
            c(1-x_{ji}a)=1.
        \]
        Equivalently,
        \[
            c-c(v^{T}u)=1.
        \]
        Now compute:
        \[
            (I_n+ucv^{T})(I_n-uv^{T})
            =
            I_n+u\bigl(c-1-c(v^{T}u)\bigr)v^{T}
            =
            I_n.
        \]
        Thus \(I_n-X(aE_{ij})\) has a left inverse for every \(X\in S\). Therefore
        \[
            aE_{ij}\in J(S).
        \]
        Since \(M_n(J(R))\) is generated additively by such elements \(aE_{ij}\), we get
        \[
            M_n(J(R))\subseteq J(S).
        \]
        
        Therefore
        \[
            J(M_n(R))=M_n(J(R)).
        \]
        \end{proof}
        
        \begin{proof}[Proof of 3]
        Let \(R\) be a PID, and let \(F\) be a finitely generated free \(R\)-module. Then
        \[
            F\cong R^n
        \]
        for some \(n\geq 0\). We prove by induction on \(n\) that every submodule of \(R^n\)
        is free.
        
        For \(n=0\), the claim is trivial. For \(n=1\), every submodule of \(R\) is an ideal of
        \(R\). Since \(R\) is a PID, every ideal is principal. Hence every submodule of \(R\)
        is free.
        
        Assume now that \(n\geq 2\), and suppose the result is true for \(R^{n-1}\). Let
        \[
            N\subseteq R^n
        \]
        be a submodule. Consider the projection onto the first coordinate:
        \[
            \pi:R^n\longrightarrow R.
        \]
        Then \(\pi(N)\) is an ideal of \(R\). Since \(R\) is a PID, there exists \(d\in R\) such that
        \[
            \pi(N)=(d).
        \]
        
        Let
        \[
            K=N\cap \ker\pi.
        \]
        Since
        \[
            \ker\pi\cong R^{n-1},
        \]
        the induction hypothesis implies that \(K\) is free.
        
        We have an exact sequence
        \[
            0\longrightarrow K\longrightarrow N\xrightarrow{\ \pi\ }\pi(N)\longrightarrow 0.
        \]
        Since
        \[
            \pi(N)=(d)
        \]
        is a principal ideal, it is a free \(R\)-module of rank \(1\) if \(d\neq 0\), and it is
        the zero module if \(d=0\). In either case, \(\pi(N)\) is projective. Hence the short
        exact sequence splits. Therefore
        \[
            N\cong K\oplus \pi(N).
        \]
        Since both \(K\) and \(\pi(N)\) are free, their direct sum is free. Hence \(N\) is free.
        
        By induction, every submodule of a finitely generated free module over a PID is free.
        \end{proof}